\chapter{物理结构设计与性能优化}

\section{数据库系统选择}

当前项目使用 openGauss 数据库，并通过 Docker 在本地运行。Docker Compose 文件位于 \code{opengauss_setup/docker/docker-compose.yml}，容器端口映射为宿主机 \code{15432} 到容器 \code{5432}。Python 端通过 \code{psycopg2} 连接，配置见 \code{config.py}。

选择 openGauss 的意义在于，它是面向企业级场景的关系数据库系统，支持标准关系模型、事务、锁、索引、约束等数据库课程核心内容，适合用于本课程大作业展示。

\section{数据类型设计}

业务编号如 \code{student_id}、\code{teacher_id}、\code{course_id} 使用 \code{VARCHAR}，便于保存带前缀或学号格式的业务编码。代理主键如 \code{user_id}、\code{offering_id}、\code{enrollment_id} 使用 \code{BIGSERIAL}，便于自动生成。成绩和学分使用 \code{DECIMAL}，避免浮点误差。日期和时间使用 \code{DATE} 与 \code{TIMESTAMP}，可支持学期范围和选课窗口判断。

\section{索引设计}

主键和唯一约束会自动生成索引。此外，当前 \code{init.sql} 显式创建了以下常用索引：

\begin{itemize}
  \item \code{idx_user_session_user} 与 \code{idx_user_session_valid}：支持持久登录会话校验。
  \item \code{idx_offering_course_semester}：支持按课程和学期查询开课班次。
  \item \code{idx_offering_teacher_semester}：支持教师按学期查看负责班次。
  \item \code{idx_schedule_offering} 与 \code{idx_schedule_time}：支持上课时间片展示和时间冲突检查。
  \item \code{idx_enrollment_student_status} 与 \code{idx_enrollment_offering_status}：支持学生成绩单、已选课程、班次人数统计和容量检查。
  \item \code{idx_score_log_enrollment} 与 \code{idx_score_log_changed_at}：支持按选课记录追踪成绩变更，以及按时间倒序查看审计日志。
  \item \code{idx_course_prereq_prereq}：支持从某门先修课反查后续课程，并辅助先修关系递归检查。
\end{itemize}

\section{规范化后的索引补偿}

规范化设计提高了数据一致性，但也会增加表连接数量。例如，课程列表需要从 \code{course_offering} 连接 \code{course}、\code{teacher}、\code{semester}、\code{classroom}，上课时间需要从 \code{course_schedule} 聚合得到，已选人数需要从 \code{enrollment} 统计得到。因此，物理设计需要通过索引补偿规范化带来的查询成本。

\code{idx_offering_course_semester} 支持按课程和学期查询教学班，适合管理员维护某门课程的开课情况；\code{idx_offering_teacher_semester} 支持教师查看指定学期教学任务；\code{idx_schedule_offering} 支持从教学班快速读取上课时间片；\code{idx_schedule_time} 支持按星期和节次进行时间冲突检测；\code{idx_enrollment_student_status} 支持学生查看自己已选、已完成或已退课程；\code{idx_enrollment_offering_status} 支持统计某教学班当前选课人数，也是容量检查的关键索引。

这组索引与第 5 章的规范化改造形成配合：系统不再在基础表中保存 \code{schedule_text}、\code{selected_count} 和 \code{gpa_point}，而是在查询时通过聚合、连接或 CASE 表达式得到展示值。索引设计的作用就是让这种“减少冗余、查询计算”的模式在数据量增长后仍具有可接受的性能。

\section{视图设计与派生数据查询}

本轮在 \code{migrate_triggers_constraints_20260608.sql} 中新增了三个查询视图，用于在不破坏 3NF 的前提下提高查询便利性。这些视图不是基础表，不承担数据持久化职责，因此不会把 \code{selected_count} 或展示用 \code{schedule_text} 重新引入基础关系模式。

\begin{longtable}{p{0.25\textwidth}p{0.35\textwidth}p{0.31\textwidth}}
\caption{规范化后的查询视图}\\
\toprule
视图 & 主要字段 & 设计目的 \\
\midrule
\endfirsthead
\toprule
视图 & 主要字段 & 设计目的 \\
\midrule
\endhead
\bottomrule
\endfoot
\code{v_offering_selected_count} & \code{offering_id}、\code{max_capacity}、\code{selected_count}、\code{remaining_capacity} & 通过聚合 \code{enrollment} 统计已选人数和剩余容量，替代基础表中的派生字段。 \\
\code{v_course_offering_detail} & 课程、教师、学期、教室、容量、已选人数、剩余容量、展示用 \code{schedule_text} & 服务课程查询、教师班次和管理员统计页面，减少应用层重复 JOIN。 \\
\code{v_student_timetable} & 学生、课程、教师、教室、星期、开始节次、结束节次 & 支持学生课表查询和时间冲突规则说明。 \\
\end{longtable}

需要强调的是，视图中的 \code{selected_count} 和 \code{schedule_text} 是查询结果，不是基础表字段。它们不会引入基础表更新异常，也不需要触发器维护。视图把规范化结构转换为页面更容易使用的展示结构，体现了物理设计对逻辑设计的补偿。

\section{高频查询}

高频查询主要包括：学生查看可选课程、学生查看已选课程、教师查看某班次学生名单、管理员查看班次容量、成绩单查询、成绩分布统计。当前服务层已经在课程列表、学生可选课程和管理员班次统计中使用 \code{v_course_offering_detail}；成绩单中的 \code{gpa_point} 仍通过 CASE 表达式按 \code{final_score} 动态计算，避免保存派生字段。

\section{性能风险与优化}

当前系统数据量较小，查询性能足以支撑课程演示。但随着数据规模扩大，动态聚合人数和时间冲突查询可能成为热点。可继续优化的方向包括：

\begin{itemize}
  \item 使用执行计划分析高频 SQL，确认索引是否被使用。
  \item 对成绩分布、班次容量等统计页面增加分页或限制。
  \item 对复杂报表考虑物化视图或定期统计表，但需权衡规范化和一致性。
  \item 对模糊查询字段增加合适索引或限制查询范围。
\end{itemize}

\section{备份恢复}

当前项目提供初始化脚本 \code{opengauss_setup/docker/init_db.sh}，可以重建数据库并导入样例数据。项目还提供 \code{migrate_3nf_20260608.sql}，用于旧结构迁移到 3NF 结构。正式运行环境中应增加定期备份、恢复演练和迁移版本管理，避免依赖手工操作。
